% =====================================================================
%  1bat-ud3-13.tex  —  HORA 13: Validació amb GeoGebra
%  (sense preàmbul: s'inclou des de main.tex)
% =====================================================================
\horatitol{Sessió 13 — Validació amb GeoGebra}%
{Comprovar amb l'ordinador el treball fet a mà: domini, punts de tall i forma de la gràfica. L'eina com a autocorrecció, no com a drecera.}

\subsection{Idea clau}

A la sessió anterior vau analitzar funcions \textbf{a mà}. Avui les portem a
\textbf{GeoGebra} per \emph{validar} aquell treball. La idea important: GeoGebra
\textbf{no fa la feina per tu}; serveix per \textbf{comprovar} si el que has calculat
és correcte. Quan coincideix, guanyes confiança; quan \emph{no} coincideix, has de
tornar al paper i \textbf{trobar l'error}.

\begin{idea}
\textbf{Com validar una funció amb GeoGebra}
\begin{enumerate}[leftmargin=1.6em, itemsep=2pt]
  \item \textbf{Introdueix la funció} a la barra d'entrada, p.\,ex. \texttt{f(x)=(x-5)/(x-2)}.
  \item \textbf{Domini:} mira on la gràfica \emph{s'interromp} (asímptotes verticals,
        forats). Hauria de coincidir amb els valors que vas excloure a mà.
  \item \textbf{Tall amb l'eix $Y$:} mira on la corba creua l'eix vertical. Pots
        escriure \texttt{f(0)} per veure'n el valor exacte.
  \item \textbf{Talls amb l'eix $X$:} usa l'eina \emph{Intersecció} o escriu
        \texttt{Arrel(f)} per obtenir on creua l'eix horitzontal.
  \item \textbf{Compara} cada resultat amb el teu càlcul. Si hi ha desajust, l'error
        és (gairebé sempre) a l'àlgebra: revisa-la.
\end{enumerate}
\end{idea}

\subsection{Exemples}

\begin{exemple}[validació correcta: una paràbola]
A mà havíem trobat que $f(x)=x^2-x-6$ té domini $\mathbb{R}$, tall amb $Y$ a
$(0,-6)$ i talls amb $X$ a $(-2,0)$ i $(3,0)$. A GeoGebra, la gràfica confirma tot:

\begin{center}
\begin{tikzpicture}
\begin{axis}[udaxis, width=9cm, height=5.4cm,
    xlabel={\small $x$}, ylabel={\small $y$},
    xmin=-4, xmax=5, ymin=-7, ymax=4,
    xtick={-2,0,3}, ytick={-6,0}]
  \addplot[udblue, domain=-3.2:4.2, samples=80]{x^2-x-6};
  \addplot[udnavy, only marks, mark=*, mark size=1.7pt] coordinates {(-2,0) (3,0) (0,-6)};
  \node[udnavy, font=\scriptsize, anchor=south east] at (axis cs:-2,0) {$(-2,0)$};
  \node[udnavy, font=\scriptsize, anchor=south west] at (axis cs:3,0) {$(3,0)$};
  \node[udnavy, font=\scriptsize, anchor=west] at (axis cs:0.15,-6) {$(0,-6)$};
\end{axis}
\end{tikzpicture}

\figcap{La gràfica de GeoGebra creua els eixos exactament als punts calculats a mà: validació correcta.}
\end{center}
\end{exemple}

\begin{exemple}[la validació detecta un error]
Suposem que algú, a mà, havia escrit que $f(x)=\dfrac{x-5}{x-2}$ talla l'eix
horitzontal a $x=2$. En portar-ho a GeoGebra, veu que la gràfica \textbf{no} creua
l'eix a $x=2$: justament allà hi ha una \textbf{asímptota vertical} (la funció no hi
està definida!). La corba creua l'eix horitzontal a $x=5$. \textbf{Conclusió:} hi
havia un error: $x=2$ és el valor \emph{exclòs del domini}, no un tall. El tall amb
l'eix $X$ és a $(5,0)$. La validació ha funcionat com a autocorrecció.
\end{exemple}

\subsection{Exercicis resolts}

\begin{exresolt}[interpretar el que mostra GeoGebra]
En introduir $f(x)=\dfrac{2x-4}{x+1}$, GeoGebra dibuixa una corba que s'enfila cap a
l'infinit a prop de $x=-1$ i que creua l'eix horitzontal en un punt. Quins valors
hauries de veure i com els llegeixes?
\solucio
\textbf{Domini:} la corba s'interromp a $x=-1$ (asímptota vertical): confirma
$\mathrm{Dom}\,f=\mathbb{R}\setminus\{-1\}$, com havíem calculat a mà (denominador
zero a $x=-1$).

\textbf{Tall amb $X$:} la corba creua l'eix horitzontal on el numerador s'anul·la,
$2x-4=0 \Rightarrow x=2$: punt $(2,0)$.

\textbf{Tall amb $Y$:} escrivint \texttt{f(0)} obtenim $\dfrac{-4}{1}=-4$: punt
$(0,-4)$.
\resposta{GeoGebra confirma domini $\mathbb{R}\setminus\{-1\}$, tall $X$ a $(2,0)$ i tall $Y$ a $(0,-4)$.}
\end{exresolt}

\begin{exresolt}[quan el càlcul i la gràfica no coincideixen]
A mà has obtingut que $g(x)=\sqrt{x-2}$ té domini $\mathbb{R}$. Però a GeoGebra la
gràfica \textbf{només apareix a la dreta de $x=2$}. Què ha passat i com ho corregeixes?
\solucio
La gràfica de GeoGebra és la pista: si la corba només existeix per a $x\ge 2$, vol
dir que el domini \textbf{no} és tot $\mathbb{R}$. L'error a mà va ser oblidar que el
\textbf{radicand ha de ser $\ge 0$}: $x-2\ge 0 \Rightarrow x\ge 2$. El domini correcte
és $[2,+\infty)$. La gràfica ha permès detectar i corregir l'error.
\resposta{Domini correcte: $[2,+\infty)$. L'error era oblidar la restricció de l'arrel;
la gràfica de GeoGebra ho ha fet evident.}
\end{exresolt}

\subsection{Exercicis per fer}

\noindent\textit{Full de pràctica tecnològica. Per a cada funció: escriu el que vas
calcular a mà (sessió 12), comprova-ho a GeoGebra i marca si coincideix o si has
trobat un error.}

\begin{exercicis}
\begin{exlist}
  \item Introdueix a GeoGebra $f(x)=x^2-9$ i comprova: el domini, el tall amb l'eix
        $Y$ i els talls amb l'eix $X$. Coincideix amb el teu càlcul a mà?
  \item Introdueix $f(x)=\dfrac{x-5}{x-2}$ i localitza l'asímptota vertical. Quin
        valor del domini queda exclòs? On talla l'eix horitzontal?
  \item Introdueix $f(x)=\sqrt{x+3}$ i observa a partir de quin valor de $x$ apareix
        la gràfica. Quin és el domini? On talla l'eix $X$?
  \item Completa el full de validació amb tres funcions a la teva elecció de la
        sessió 12:
        \begin{center}
        \renewcommand{\arraystretch}{1.4}
        \begin{tabular}{p{3.2cm} p{3.2cm} p{2.2cm} p{2.6cm}}
        \toprule
        \textbf{Funció} & \textbf{El que vaig calcular a mà} & \textbf{Coincideix?} & \textbf{Error trobat (si n'hi ha)} \\
        \midrule
         & & & \\[10pt]
         & & & \\[10pt]
         & & & \\
        \bottomrule
        \end{tabular}
        \end{center}
  \item \textbf{(Autocorrecció)} Si en alguna funció el resultat de GeoGebra
        \emph{no} coincideix amb el teu, torna al paper, troba l'error i escriu en una
        frase quin va ser (de domini? de tall? de signe en resoldre l'equació?).
  \item \textbf{(Coherència)} Explica, amb les teves paraules, per què una asímptota
        vertical a GeoGebra correspon sempre a un valor \textbf{exclòs del domini} i
        \textbf{mai} a un tall amb l'eix horitzontal.
\end{exlist}
\end{exercicis}

% ---------------------------------------------------------------------
\fullsolucions{Sessió 13}
\begin{solucions}
\begin{sollist}
  \item $\mathrm{Dom}=\mathbb{R}$; tall $Y$: $(0,-9)$; talls $X$: $(3,0)$ i $(-3,0)$.
        La gràfica (paràbola cap amunt) ho confirma.
  \item Asímptota vertical a $x=2$: queda exclòs el $2$
        ($\mathrm{Dom}=\mathbb{R}\setminus\{2\}$). Talla l'eix horitzontal a $(5,0)$
        (numerador zero).
  \item La gràfica apareix a partir de $x=-3$: $\mathrm{Dom}=[-3,+\infty)$. Talla
        l'eix $X$ a $(-3,0)$.
  \item Resposta oberta (full de validació personal). Es valora que el càlcul a mà i
        la lectura de GeoGebra siguin coherents i que els desajustos quedin anotats.
  \item Resposta oberta. Es valora identificar correctament el tipus d'error
        (domini, tall o signe) i corregir-lo.
  \item Una asímptota vertical apareix on la funció \textbf{no està definida} (p.\,ex.
        denominador zero): és un valor que \emph{no} pertany al domini, i per tant no
        pot ser-hi cap punt de la gràfica, ni un tall. Un tall amb l'eix horitzontal, en
        canvi, és un punt \emph{real} de la gràfica on $y=0$; són coses oposades.
\end{sollist}
\end{solucions}
